\begin{longtable}{|P{0.15\linewidth}|P{0.17\linewidth}|P{0.22\linewidth}|P{0.43\linewidth}|}
\caption{STBS IoT network connectivity.}\label{tab:a1b-network-connectivity}\\
\hline
\rowcolor{gray!20}\textbf{IoT devices, gateways, routers, etc. to connect} & \textbf{Main considerations} & \textbf{Specific questions} & \textbf{Responses}\\
\hline
\endfirsthead
\hline
\rowcolor{gray!20}\textbf{IoT devices, gateways, routers, etc. to connect} & \textbf{Main considerations} & \textbf{Specific questions} & \textbf{Responses}\\
\hline
\endhead
\hline
\multicolumn{4}{r}{\footnotesize Continued on next page}\\
\endfoot
\hline
\endlastfoot
Selected devices: ESP32-C6 \#1/\#2; field sensors and actuators; Aqara C1 feeder; Reolink camera; Zigbee coordinator; Raspberry Pi 5; monitor; switch; router. & Determine the connectivity between selected IoT devices, gateways, routers, etc. & Why do selected devices need to be connected with each other and to the network? & The water-environment detector connects DFR0198, SEN0711 and SEN0189 to ESP32-C6 \#1. The tank-functionality detector connects SEN0257 and SEN0217 to ESP32-C6 \#2 and drives the valve and AC electric water-cycling pump. The two ESP32-C6 nodes and Aqara C1 food feeder use the Zigbee path to the coordinator, while the Reolink camera uses Ethernet. The Raspberry Pi 5 aggregates local status, camera data and feeding events before forwarding selected information to the cloud.\\
\cline{2-4}
 & Determine communication and network technologies required to support connectivity & What are the range and coverage requirements? & DFR0198 uses 1-Wire/GPIO; SEN0711 uses RS-485 Modbus through the DFR0845 isolated converter to UART; SEN0189 and SEN0257 use analog GPIO/ADC with 3.3 V conditioning; and SEN0217 uses conditioned Hall pulses on GPIO. Zigbee links use 2.4\,GHz IEEE 802.15.4 and are shown as no more than 10\,m from the ESP32-C6 nodes or Aqara feeder to the coordinator. The Reolink camera, Raspberry Pi 5, LAN switch and router use Ethernet over RJ45; HDMI connects the local monitor and USB connects the Zigbee dongle.\\
\cline{3-4}
 &  & What are the power consumption constraints? & The ESP32-C6 controllers, sensors, Zigbee coordinator and Aqara feeder require low-voltage local supplies; the valve uses a 24 VDC supply; the heater and water-cycling pump are mains-powered electrical actuators; and the Reolink camera can use PoE. The pump is not battery-powered. The LAN switch, Raspberry Pi 5, monitor and router require suitable local power. RCD/overload protection and a local backup should preserve the gateway, coordinator and controllers long enough to record alarms and keep bounded local control available.\\
\cline{3-4}
 &  & What is the required data rate (how much data needs to be transmitted over time)? & Temperature, pH/NH3-N, turbidity, level and flow readings, actuator states and Aqara feeder events are low-rate telemetry. The IP camera is the bandwidth-heavy device and therefore uses Ethernet to the LAN switch. The Raspberry Pi 5 aggregates and filters the local streams, so the cloud receives selected telemetry, events or processed video rather than requiring every raw local signal to travel continuously over the Internet.\\
\cline{3-4}
 &  & What are the latency requirements (how fast does data need to be transmitted)? & Heater, valve and AC pump decisions use local relay, PWM or driver signals and do not wait for the cloud. Sensor feedback reaches the ESP32-C6 controllers over 1-Wire/GPIO, RS-485/UART, GPIO/ADC or Pulse/GPIO links. Zigbee and Ethernet support gateway monitoring and event delivery, while MQTT/HTTPS cloud communication is used for reporting and alerts rather than immediate safety action.\\
\cline{3-4}
 &  & What are the security and reliability requirements? & Local control should continue when the router or Internet is unavailable. The Raspberry Pi 5 stores data locally when upload is interrupted and the gateway labels or forwards buffered events after reconnection. MQTT/HTTPS traffic to the cloud uses TLS; Zigbee devices require authenticated joining and coordinator access control; and actuator commands require local bounds and interlocks. RS-485 isolation, ADC protection, GPIO level conditioning, RCD/overload protection, cable glands and waterproof enclosures are required.\\
\cline{3-4}
 &  & What are the cost considerations? & ESP32-C6 nodes and short Zigbee links reduce point-to-point wiring inside the device layer. Ethernet is reserved for the Reolink IP camera and the fixed fog/network path, while one Raspberry Pi 5 provides local aggregation, buffering and display output. The selected sensors require signal conditioning and protected enclosures; the AC pump and heater require correctly rated switching and electrical protection. Cloud storage and analytics can be scaled according to selected uploads.\\
\cline{3-4}
 &  & Which communication/network technologies are selected? & The diagram selects 1-Wire/GPIO for temperature; RS-485 Modbus through DFR0845 and UART for the combined pH/NH3-N sensor; GPIO/ADC for turbidity and pressure-derived level; Pulse/GPIO for flow; Relay/PWM/Driver for heater, valve and AC pump control; Zigbee at 2.4\,GHz under IEEE 802.15.4 for the ESP32-C6 and Aqara feeder links; Ethernet at 10/100/1000 Mbps over RJ45 for the Reolink camera and fog/network path; USB for the coordinator; HDMI for the local monitor; and MQTT/HTTPS over TLS for the Internet connection to the IoT cloud platform.\\
\end{longtable}
